\section{The Fear of Death Undone}

The Son of God did not abolish death by forbidding it to touch him. He entered it. ``Therefore because the children are partakers of flesh and blood, he also himself in like manner hath been partaker of the same: that, through death, he might destroy him who had the empire of death, that is to say, the devil: and might deliver them, who through the fear of death were all their lifetime subject to servitude.''\endnote{Hebrews 2:14--15, Douay--Rheims (Challoner).} The Apostle names the bondage exactly. It is not death that enslaves, since death comes only once; it is the fear of death, which lasts a whole lifetime and bends every choice toward the one thing it cannot avoid. Christ breaks that servitude from within. He took the death that belongs to the human condition and accepted it in free submission to his Father's will, and the Church confesses the result in one sentence: ``The obedience of Jesus has transformed the curse of death into a blessing.''\endnote{CCC 1009.}

\subsection{Die, that thou sin not}

Augustine asks why, if death is the punishment of sin, the baptized still die, since their guilt has been washed away. His answer goes to the root of the matter. God left the parting of soul and body in place ``for this reason, that if the immortality of the body followed immediately upon the sacrament of regeneration, faith itself would be thereby enervated. For faith is then only faith when it waits in hope for what is not yet seen in substance.''\endnote{Augustine, \emph{De civitate Dei} XIII.4, Dods.} If baptism ended death, men would run to the font to escape the grave and not to find Christ, and there would be no martyrs, since there would be nothing to suffer. Instead, ``by the greater and more admirable grace of the Saviour, the punishment of sin is turned to the service of righteousness. For then it was proclaimed to man, `If thou sinnest, thou shalt die;' now it is said to the martyr, `Die, that thou sin not.'\,''\endnote{\emph{De civitate Dei} XIII.4, Dods.}

The threat that once guarded the commandment now keeps it. Augustine presses the inversion until it is complete: ``For then death was incurred by sinning, now righteousness is fulfilled by dying.'' Yet he will not let anyone conclude that death has become good. ``Not that death, which was before an evil, has become something good, but only that God has granted to faith this grace, that death, which is the admitted opposite to life, should become the instrument by which life is reached.''\endnote{\emph{De civitate Dei} XIII.4, Dods.} Death has not been pardoned. It has been conscripted. The enemy remains an enemy, and serves.

That is why the Christian need not choose between dread and indifference. He fears death rightly, as a punishment and a wrench, and he hopes through it rightly, because Christ has passed through it before him and made it the road home. Augustine says in the next chapter that the good ``die well, though death is an evil.''\endnote{\emph{De civitate Dei} XIII.5, Dods.}

\subsection{The martyrs' longing}

The first Christians who faced death for Christ show what this looks like in a human face. Ignatius of Antioch, carried under guard to Rome to be thrown to the beasts, wrote ahead to the Christians of that city and begged them not to obtain his release. ``All the pleasures of the world, and all the kingdoms of this earth, shall profit me nothing. It is better for me to die in behalf of Jesus Christ, than to reign over all the ends of the earth. \ldots{} Him I seek, who died for us: Him I desire, who rose again for our sake.'' And he turns their kindness inside out: ``do not hinder me from living, do not wish to keep me in a state of death.''\endnote{Ignatius of Antioch, \emph{Epistle to the Romans} 6 (shorter recension), trans. A. Roberts and J. Donaldson, \emph{Ante-Nicene Fathers}, vol.~1.} For the martyr, to go on living in this world apart from Christ is the real death, and the arena is the door to life.

Polycarp of Smyrna was an old man when the proconsul urged him to curse Christ and go free. ``Eighty and six years have I served Him, and He never did me any injury: how then can I blaspheme my King and my Saviour?''\endnote{\emph{The Martyrdom of Polycarp} 9, \emph{Ante-Nicene Fathers}, vol.~1.} Bound at the stake, he looked up to heaven and gave thanks: ``I give Thee thanks that Thou hast counted me worthy of this day and this hour, that I should have a part in the number of Thy martyrs, in the cup of thy Christ, to the resurrection of eternal life, both of soul and body.''\endnote{\emph{Martyrdom of Polycarp} 14, \emph{Ante-Nicene Fathers}, vol.~1.} His church gathered his bones afterward, ``as being more precious than the most exquisite jewels,'' and resolved to keep each year the anniversary of his martyrdom, which their Greek calls his birthday.\endnote{\emph{Martyrdom of Polycarp} 18, \emph{Ante-Nicene Fathers}, vol.~1, with the translators' note: ``Literally, `the birth-day.'\,''} The Church still keeps the feasts of her saints, for the most part, on the day of their death. In her calendar a Christian is born when he dies.

Few are asked to die as Ignatius and Polycarp died. But every Christian is asked to die in Christ, and the Apostle gives the pattern for all of them: ``For to me, to live is Christ: and to die is gain. And if to live in the flesh: this is to me the fruit of labour. And what I shall choose I know not. But I am straitened between two: having a desire to be dissolved and to be with Christ, a thing by far the better.''\endnote{Philippians 1:21--23, Douay--Rheims (Challoner).} Paul does not love death. He loves Christ, and he finds that death, which was the wages of sin, has become the way to him.
